\section*{\texorpdfstring{\S\CurrentExerciseGroup}{Section \CurrentExerciseGroup}}
\phantomsection
\addcontentsline{toc}{section}{\texorpdfstring{Exercises for \S\CurrentExerciseGroup}{Exercises for Section \CurrentExerciseGroup}}

¶ 1) a) Let X be a locally compact space. Show that every compact convex subset A of the locally convex space \(\mathcal{H}(\mathrm{X};\mathbf{C})\) that is contained in \(\mathcal{H}_{+}(\mathrm{X})\) is strictly compact. (Consider the upper envelope function of the set of all \(f\in\mathrm{A}\) , \(s=\sup_{f\in\mathrm{A}}f\) , which is

continuous on X and tends to 0 at the point at infinity of X. For every continuous numerical function g on X, denote by \(\mathrm{P}(g)\) the set of \(x \in X\) such that \(g(x) > 0\) . Show that there exists a function f in A such that \(\mathrm{P}(f) = \mathrm{P}(s)\) ; for this, observe that if \(K_{n}\) is the compact set of \(x \in X\) such that \(s(x) \geqslant 1/n\) , there exists a finite number of functions \(f_{n,i} \in A\) such that the open sets \(\mathrm{P}(f_{n,i})\) cover \(K_{n}\) . Then make use of the fact that A is convex and closed.)

\begin{enumerate}
\def\labelenumi{\alph{enumi})}
\setcounter{enumi}{1}
\tightlist
\item
  Assume that there exists in X a countable dense subset D. Show that every compact convex subset A of \(\mathcal{K}(X; \mathbf{C})\) is then strictly compact. (Let U be the set of \(x \in X\) such that there exists an \(f \in A\) for which \(f(x) \neq 0\) , which is an open set in X. Show, making use of the fact that A is a Baire space, that there exists a function \(g \in A\) such that \(g(x) \neq 0\) for all \(x \in U \cap D\) .)
\end{enumerate}

\begin{enumerate}
\def\labelenumi{\arabic{enumi})}
\setcounter{enumi}{1}
\tightlist
\item
  Let Y be a noncompact locally compact space that is countable at infinity and let \((Y_{n})\) be an increasing sequence of compact subsets of Y forming a covering of Y and such that \(Y_{n} \subset Y_{n+1}\) for all n, so that every compact subset of Y is contained in one of the \(Y_{n}\) (GT, I, §9, No.~9, Cor. 1 of Prop. 15). Let \(\beta Y\) be the Stone--Čech compactification of Y (GT, IX, §1, Exer. 7), so that the space \(\mathcal{C}(\beta Y; \mathbf{R})\) may be identified with the space of bounded continuous real-valued functions on Y, and Y is a dense open set in \(\beta Y\) ; for every function \(f \in \mathcal{C}(\beta Y; \mathbf{R})\) , \(\operatorname{Supp}(f)\) is thus the closure in \(\beta Y\) of the set of \(y \in Y\) such that \(f(y) \neq 0\) . Recall that if \(Z_{1}, Z_{2}\) are two subsets of Y, closed in Y and disjoint, then their closures \(\overline{Z}_{1}, \overline{Z}_{2}\) in \(\beta Y\) are also disjoint (GT, IX, §4, Exer. 17).
\end{enumerate}

Let \(\omega\) be a point of \(\beta\mathrm{Y}-\mathrm{Y}\) and let \(\mathrm{X}=\beta\mathrm{Y}-\{\omega\}\) . Show that, on the space \(\mathcal{K}(\mathrm{X};\mathbf{R})\) , the direct limit topology \(\mathcal{T}\) defined in No.~1 is identical to the topology of uniform convergence. (If \(p\) is a semi-norm that is continuous for \(\mathcal{T}\) , show, by means of a contradiction argument, that there exists a number \(c > 0\) such that \(p(f) \leqslant c \| f \|\) for every function \(f \in \mathcal{K}(\mathrm{X};\mathbf{R})\) . In the contrary case, there would exist a sequence \((f_n)\) of functions in \(\mathcal{K}(\mathrm{X};\mathbf{R})\) such that \(\| f_n \| \leqslant 1\) , \(p(f_n) \geqslant n\) , \(\operatorname{Supp}(f_n) \cap Y_n = \emptyset\) , and

\[
\operatorname{Supp} (f _ {n}) \cap \operatorname{Supp} (f _ {k}) = \varnothing
\]

for \(k < n\) . Then let \(Z_{i}\) ( \(i = 0,1\) ) be the union of the sets \(Y \cap \operatorname{Supp}(f_{2n + i})\) ; show that one of the two sets \(Z_{0}, Z_{1}\) is relatively compact in X (GT, X, §4, Exer. 7), and from this deduce a contradiction.)

From this, deduce that the space \(\mathcal{K}(\mathbf{X};\mathbf{R})\) equipped with the topology \(\mathcal{T}\) is not quasi-complete.

\begin{enumerate}
\def\labelenumi{\arabic{enumi})}
\setcounter{enumi}{2}
\tightlist
\item
  In the example of Exer. 2, take for Y the discrete space N. For every \(n \in N\) , let \(e_n\) be the function in \(\mathcal{K}(X; \mathbf{R})\) such that \(e_n(n) = 1\) and \(e_n(x) = 0\) for \(x \neq n\) in X. Show that the set A formed by the functions 0 and \(e_n/(n + 1)\) is compact in \(\mathcal{K}(X; \mathbf{R})\) , but not strictly compact. If U is the ultrafilter on N induced by the filter of neighborhoods of \(\omega\) in \(\beta Y = \beta N\) , show that, with respect to U, the family of positive measures ( \(n\varepsilon_n\) ) tends to 0 uniformly in every strictly compact subset of \(\mathcal{K}(X; \mathbf{R})\) , but does not tend uniformly to 0 in A.
\end{enumerate}

¶ 4) a) Let \(\mathbf{B}_n\) be the closed unit ball in the Euclidean space \(\mathbf{R}^n\) , and let \(\mathbf{A}_n\) be the set of functions \(f \in \mathcal{C}(\mathbf{B}_n; \mathbf{R})\) of the form

\[
f (x) = \lambda (\| x \| ^ {2} - 1) + (x | a)
\]

with \(|\lambda| \leqslant 1/n\) and \(\|a\| \leqslant 1/n\) . The set \(A_n\) is convex and compact in the Banach space \(\mathcal{C}(\mathbf{B}_n; \mathbf{R})\) , and possesses the following property: for any functions \(f_1, \ldots, f_n\) belonging to \(A_n\) , there exists a point \(x \in B_n\) such that \(f_1(x) = \ldots = f_n(x) = 0\) ; however, there exist \(n + 1\) functions \(g_1, \ldots, g_{n+1}\) belonging to \(A_n\) that are not simultaneously zero at

any point of \(\mathbf{B}_n\) (one may consider the mapping \(x \mapsto x / (1 - \|x\|^2)\) of \(\mathbf{B}_n\) into \(\mathbf{R}^n\) ). b) Let \(\mathbf{B}_n'\) be the set \(\mathbf{B}_n\) equipped with the discrete topology, and let \(\mathrm{D}_n = \beta \mathrm{B}_n'\) . Let Y (resp. \(\mathrm{Y}'\) ) be the locally compact topological sum space of the spaces \(\mathrm{D}_n\) (resp. \(\mathrm{B}_n'\) ). Show that \(\beta \mathrm{Y} = \beta \mathrm{Y}'\) .

\begin{enumerate}
\def\labelenumi{\alph{enumi})}
\setcounter{enumi}{2}
\tightlist
\item
  For every integer \(n\) , let \(\mathrm{K}_n \subset \mathcal{C}(\mathrm{D}_n; \mathbf{R})\) be the set of all real-valued functions whose restriction to \(\mathrm{B}_n' = \mathbf{B}_n\) belongs to \(\mathrm{A}_n\) . Let \(\mathrm{K}\) be the subset of \(\mathcal{C}(\beta\mathrm{Y}; \mathbf{R})\) formed by the functions whose restriction to \(\mathrm{D}_n\) belongs to \(\mathrm{K}_n\) for every \(n \in \mathbf{N}\) . Show that \(\mathrm{K}\) is a compact convex subset of the Banach space \(\mathcal{C}(\beta\mathrm{Y}; \mathbf{R})\) . For every finite family \((f_i)_{1 \leqslant i \leqslant n}\) of functions in \(\mathrm{K}\) , let \(W(f_1, \ldots, f_n)\) be the set of \(y \in Y'\) such that \(f_i(y) = 0\) for \(1 \leqslant i \leqslant n\) . Show that the \(W(f_1, \ldots, f_n)\) form a base of a filter \(\mathfrak{W}\) on the discrete space \(Y'\) . From this, conclude that there exists a point \(\omega \in \beta\mathrm{Y} - \mathrm{Y}\) such that the trace on \(Y'\) of the filter of neighborhoods of \(\omega\) in \(\beta\mathrm{Y}\) is a filter finer than \(\mathfrak{W}\) . Setting \(X = \beta Y - \{\omega\}\) , deduce from this that \(K \subset \mathcal{K}(X; \mathbf{R})\) and, with the help of Exercise 2, conclude that in \(\mathcal{K}(X; \mathbf{R})\) the set \(K\) is convex and compact, but not strictly compact.
\end{enumerate}

\begin{enumerate}
\def\labelenumi{\arabic{enumi})}
\setcounter{enumi}{4}
\tightlist
\item
  Show that if \(\mathbf{E}\) is a Fréchet space and \(\mathbf{X}\) is a locally compact space, then the space \(\mathcal{H}(\mathbf{X};\mathbf{E})\) is bornological.
\end{enumerate}

From this, deduce that for every set S of bounded subsets of \(\mathcal{K}(X;C)\) that is a covering of \(\mathcal{K}(X;C)\) and is such that every set formed by the points of a convergent sequence in \(\mathcal{K}(X;C)\) belongs to S, the space \(\mathcal{M}(X_{\mathfrak{S}};C)\) is complete.

\begin{enumerate}
\def\labelenumi{\arabic{enumi})}
\setcounter{enumi}{5}
\tightlist
\item
  Let \(\mu\) be a measure on a locally compact space \(X\) . Show that \(|\mu| = \sup \mathcal{R}(\alpha \mu)\) , where the scalar \(\alpha\) runs over the set of complex numbers of absolute value \(\leqslant 1\) . (Observe that for every function \(g \in \mathcal{K}(X; \mathbf{C})\) and every \(\varepsilon > 0\) , there exist a finite number of mappings \(h_i\) of \(X\) into \([0,1]\) , with compact supports, and scalars \(\alpha_i\) such that \(\| g - \sum_{i} \alpha_i h_i \| \leqslant \varepsilon\) and \(\| |g| - \sum_{i} |\alpha_i| h_i \| \leqslant \varepsilon\) .)
\end{enumerate}

Show that also \(|\mu| = \sup \mathcal{R}(h \cdot \mu)\) , where \(h\) runs over the set of functions in \(\mathcal{K}(\mathrm{X};\mathbf{C})\) such that \(\| h\| \leqslant 1\) .

\begin{enumerate}
\def\labelenumi{\arabic{enumi})}
\setcounter{enumi}{6}
\item
  Give an example showing that in \(\mathcal{M}(\mathrm{X};\mathbf{C})\) , the set of bounded real measures \(\mu\) such that \(\|\mu\|=a\) is not necessarily vaguely closed, even if X is compact; show that if X is locally compact and noncompact, then the set of bounded positive measures such that \(\|\mu\|=1\) is not vaguely closed.
\item
  Let \(\mathbf{X}\) be a locally compact space; for every compact subset \(\mathbf{K}\) of \(\mathbf{X}\) and every integer \(n > 0\) , let \(S_{\mathbf{K},n}\) be the set of functions \(f \in \mathcal{K}(\mathbf{X};\mathbf{C})\) with support contained in \(\mathbf{K}\) and such that \(\| f\| \leqslant n\) . One calls quasi-strong topology on \(\mathcal{M}(\mathbf{X};\mathbf{C})\) the topology of uniform convergence in the sets \(S_{\mathbf{K},n}\) . It is coarser than the strong topology (when \(\mathcal{M}(\mathbf{X};\mathbf{C})\) is regarded as the dual of the space \(\mathcal{K}(\mathbf{X};\mathbf{C})\) , cf.~TVS, III, §3, No.~1) and is identical to it when \(\mathbf{X}\) is paracompact.
\end{enumerate}

\begin{enumerate}
\def\labelenumi{\alph{enumi})}
\item
  Show that the quasi-strong topology on \(\mathcal{M}(\mathrm{X};\mathbf{C})\) is defined by the semi-norms \(\mu \mapsto |\mu |(f)\) , where \(f\) runs over \(\mathcal{K}_{+}(\mathrm{X})\) ; the space \(\mathcal{M}(\mathrm{X};\mathbf{C})\) is complete for the quasi-strong topology. The bounded subsets of \(\mathcal{M}(\mathrm{X};\mathbf{C})\) for the quasi-strong topology are the equicontinuous subsets of \(\mathcal{M}(\mathrm{X};\mathbf{C})\) .
\item
  Show that on \(\mathcal{M}(\mathrm{X};\mathbf{R})\) , the quasi-strong topology is compatible with the ordered vector space structure (TVS, II, §2, No.~7). From this, deduce that if H is an increasing directed set in \(\mathcal{M}(\mathrm{X};\mathbf{R})\) that is bounded above, then the supremum of H in \(\mathcal{M}(\mathrm{X};\mathbf{R})\) is identical to the limit, for the quasi-strong topology, of its section filter (loc. cit.).
\end{enumerate}

\begin{enumerate}
\def\labelenumi{\arabic{enumi})}
\setcounter{enumi}{8}
\tightlist
\item
  Let \(\mathbf{X}\) be the compact interval {[}0,1{]} of \(\mathbf{R}\) .
\end{enumerate}

\begin{enumerate}
\def\labelenumi{\alph{enumi})}
\item
  Let \(\mu_{n}\) be the measure defined by the mass \(+1\) at the point 0 and the mass \(-1\) at the point \(1/n\) . Show that the sequence \((\mu_{n})\) tends vaguely to 0 in \(\mathcal{M}(\mathrm{X};\mathbf{C})\) , but that \(|\mu_n|\) does not tend vaguely to 0.
\item
  Let \(\mu\) be Lebesgue measure on X and let \(g_{n}(x) = \sin nx\) . Show that the measure \((1 - g_{n}) \cdot \mu\) is positive and tends vaguely to \(\mu\) as \(n\) tends to infinity, but does not tend strongly to \(\mu\) . From this, deduce that the topologies induced on \(\mathcal{M}_{+}(X)\) by the vague topology and the strong topology (which is identical here to the quasi-strong topology) are distinct (cf.~Exer. 10).
\end{enumerate}

\begin{enumerate}
\def\labelenumi{\arabic{enumi})}
\setcounter{enumi}{9}
\tightlist
\item
  Let \(\Phi\) be a filter on the set \(\mathcal{M}_{+}(\mathrm{X})\) of positive measures on \(\mathrm{X}\) . Show that if \(\Phi\) is vaguely convergent then, for every compact subset \(\mathrm{K}\) of \(\mathrm{X}\) , there exist a set \(\mathrm{M} \in \Phi\) and a number \(a_{\mathrm{K}} > 0\) such that \(|\mu(f)| \leqslant a_{\mathrm{K}} \cdot \|f\|\) for every function \(f \in \mathcal{K}(\mathrm{X}, \mathrm{K}; \mathbf{C})\) and every measure \(\mu \in \mathrm{M}\) .
\end{enumerate}

¶ 11) a) Show that when \(\mathcal{M}(\mathrm{X};\mathbf{C})\) is equipped with the quasi-strong topology (resp. the topology of strictly compact convergence, resp. the vague topology), the bilinear form \((f,\mu)\mapsto\langle f,\mu\rangle\) is \((\mathfrak{S},\mathfrak{T})\) -hypocontinuous, where \(\mathfrak{T}\) is the set of vaguely bounded subsets of \(\mathcal{M}(\mathrm{X};\mathbf{C})\) , and \(\mathfrak{S}\) is the set of bounded subsets of \(\mathcal{H}(\mathrm{X};\mathbf{C})\) that are contained in \(\mathcal{H}(\mathrm{X},\mathrm{K};\mathbf{C})\) for a suitable compact set K (resp. the set of strictly compact subsets of \(\mathcal{H}(\mathrm{X};\mathbf{C})\) , resp. the set of finite subsets of \(\mathcal{H}(\mathrm{X};\mathbf{C})\) ) (cf.~TVS, III, §5, Exer. 12).

\begin{enumerate}
\def\labelenumi{\alph{enumi})}
\setcounter{enumi}{1}
\item
  For every compact subset \(\mathbf{K}\) of \(\mathbf{X}\) , show that the bilinear form \((f, \mu) \mapsto \langle f, \mu \rangle\) is continuous on \(\mathcal{H}(\mathbf{X}, \mathbf{K}; \mathbf{C}) \times \mathcal{M}(\mathbf{X}; \mathbf{C})\) when \(\mathcal{M}(\mathbf{X}; \mathbf{C})\) is equipped with the quasi-strong topology; it is continuous on \(\mathcal{H}(\mathbf{X}, \mathbf{K}; \mathbf{C}) \times \mathcal{M}_{+}(\mathbf{X})\) when \(\mathcal{M}_{+}(\mathbf{X})\) is equipped with the vague topology (make use of Exer. 10).
\item
  Assume that \(X\) is paracompact and noncompact. Show that the bilinear form \((f, \mu) \mapsto \langle f, \mu \rangle\) is not continuous on \(\mathcal{K}(X; \mathbf{C}) \times \mathcal{M}_{+}(X)\) when \(\mathcal{M}_{+}(X)\) is equipped with the strong topology (identical in this case to the quasi-strong topology).
\item
  Give an example where X is compact and the bilinear form \((f,\mu)\mapsto\langle f,\mu\rangle\) is not continuous on \(\mathcal{C}(\mathrm{X};\mathbf{C})\times\mathcal{M}(\mathrm{X};\mathbf{C})\) when \(\mathcal{M}(\mathrm{X};\mathbf{C})\) is equipped with the topology of compact convergence (TVS, IV, §1, Exer. 11).
\end{enumerate}

¶ 12) a) Show that the bilinear mapping \((g, \mu) \mapsto g \cdot \mu\) of \(\mathcal{C}(\mathrm{X}; \mathbf{C}) \times \mathcal{M}(\mathrm{X}; \mathbf{C})\) into \(\mathcal{M}(\mathrm{X}; \mathbf{C})\) is continuous when \(\mathcal{C}(\mathrm{X}; \mathbf{C})\) is equipped with the topology of compact convergence and \(\mathcal{M}(\mathrm{X};\mathbf{C})\) with the quasi-strong topology (Exer. 8) or with the topology of strictly compact convergence.

\begin{enumerate}
\def\labelenumi{\alph{enumi})}
\setcounter{enumi}{1}
\item
  Show that the bilinear mapping \((g,\mu)\mapsto g\cdot \mu\) of \(\mathcal{C}(\mathrm{X};\mathbf{C})\times \mathcal{M}(\mathrm{X};\mathbf{C})\) into \(\mathcal{M}(\mathrm{X};\mathbf{C})\) is \((\mathfrak{S},\mathfrak{T})\) -hypocontinuous when \(\mathcal{C}(\mathrm{X};\mathbf{C})\) is equipped with the topology of compact convergence and \(\mathcal{M}(\mathrm{X};\mathbf{C})\) with the vague topology, where \(\mathfrak{S}\) denotes the set of finite subsets of \(\mathcal{C}(\mathrm{X};\mathbf{C})\) and \(\mathfrak{T}\) the set of vaguely bounded subsets of \(\mathcal{M}(\mathrm{X};\mathbf{C})\) .
\item
  Show that the mapping \((g, \mu) \mapsto g \cdot \mu\) of \(\mathcal{C}(\mathrm{X};\mathbf{C}) \times \mathcal{M}_{+}(\mathrm{X})\) into \(\mathcal{M}(\mathrm{X};\mathbf{C})\) is continuous when \(\mathcal{C}(\mathrm{X};\mathbf{C})\) is equipped with the topology of compact convergence, \(\mathcal{M}(\mathrm{X};\mathbf{C})\) and \(\mathcal{M}_{+}(\mathrm{X})\) with the vague topology.
\item
  Give an example of a compact space X such that the bilinear mapping \((g,\mu)\mapsto g\cdot\mu\) of \(\mathcal{C}(\mathrm{X};\mathbf{C})\times\mathcal{M}(\mathrm{X};\mathbf{C})\) into \(\mathcal{M}(\mathrm{X};\mathbf{C})\) is not continuous when \(\mathcal{C}(\mathrm{X};\mathbf{C})\) is equipped with the topology of uniform convergence and \(\mathcal{M}(\mathrm{X};\mathbf{C})\) with the vague topology.
\end{enumerate}

\begin{enumerate}
\def\labelenumi{\arabic{enumi})}
\setcounter{enumi}{12}
\tightlist
\item
  Show that for every function \(f \in \mathcal{K}_{+}(X)\) , the mapping \(\mu \mapsto |\mu|(f)\) is lower semi-continuous on \(\mathcal{M}(X; \mathbf{C})\) for the vague topology.
\end{enumerate}

¶ 14) a) Let X be a locally compact space with a countable base. Show that the space \(\mathcal{M}_{+}(X)\) , equipped with the vague topology, is metrizable.

\begin{enumerate}
\def\labelenumi{\alph{enumi})}
\setcounter{enumi}{1}
\tightlist
\item
  Let L be the set of increasing sequences \((k_{n})_{n\in\mathbf{Z}}\) of integers \(\geqslant0\) , tending to \(+\infty\) with n, and zero for \(n\leqslant0\) except for a finite number of values of n; let P be the set of equivalence classes formed by the translates \((k_{n+h})_{n\in\mathbf{Z}}\) of a same sequence \((k_{n})_{n\in\mathbf{Z}}\) (where h runs over Z); L and P are uncountable. Let X be the discrete topological sum space of P and Z; let H be the set of positive measures on X defined in the following way: for every \(\alpha\in P\) and every sequence \(g\in\alpha\) , \(\nu_{\alpha,g}\) is the measure defined by the mass +1 at the point \(\alpha\) , mass 0 at the other points of P, and mass \(g(n)\) at every point \(n\in\mathbf{Z}\) ; take H to be the cone generated by the measures \(\nu_{\alpha,g}\) . Now let \(\mu\) be the measure on X defined by the mass +1 at each point of P and mass 0 at each point of Z. Show that \(\mu\) is in the vague closure of H but is not in the vague closure of any vaguely bounded subset of H. (Make use of the fact that for every mapping f of Z into N, there exists an increasing sequence \((k_{n})\in L\) such that for every \(h\in Z\) , \(\lim_{n\to\infty}k_{n+h}/f(n)=+\infty\) .)
\end{enumerate}

\begin{enumerate}
\def\labelenumi{\arabic{enumi})}
\setcounter{enumi}{14}
\tightlist
\item
  Let \(\mathbf{X}\) be a noncompact locally compact space; on the space \(\mathcal{M}^1(\mathbf{X};\mathbf{C})\) of bounded measures on \(\mathbf{X}\) , the dual of the Banach space \(\mathcal{C}^0(\mathbf{X};\mathbf{C})\) , the topology defined by the norm \(\| \mu \|\) is called the ultrastrong topology, and the topology \(\sigma (\mathcal{M}^1(\mathbf{X};\mathbf{C}),\mathcal{C}^0(\mathbf{X};\mathbf{C}))\) the weak topology.
\end{enumerate}

\begin{enumerate}
\def\labelenumi{\alph{enumi})}
\item
  Show that if \(X\) is paracompact, then the weak topology on \(\mathcal{M}^1(X; \mathbf{C})\) is strictly finer than the vague topology (compare with Exer. 2).
\item
  Show that on \(\mathcal{M}^1 (\overline{\mathrm{X}};\mathbf{C})\) , the ultrastrong topology is strictly finer than the quasi-strong topology (observe that for the latter, \(\varepsilon_{x}\) tends to 0 as \(x\) tends to the point at infinity of X).
\item
  If \(\mathbf{X}\) is countable at infinity and is not discrete, show that on \(\mathcal{M}^1(\mathbf{X};\mathbf{C})\) the weak topology and the quasi-strong topology are not comparable.
\item
  In \(\mathcal{M}^1 (\mathrm{X};\mathbf{C})\) , every weakly bounded set is bounded for the ultrastrong topology, but there can exist sets that are bounded for the quasi-strong topology but are not weakly bounded.
\end{enumerate}

\begin{enumerate}
\def\labelenumi{\arabic{enumi})}
\setcounter{enumi}{15}
\item
  Let X be a locally compact space, \(X'\) the compact space obtained by adjoining to X a point at infinity. Let \(\mu\) be a bounded measure on X; show that there exists one and only one measure \(\mu'\) on \(X'\) that extends \(\mu\) and satisfies \(\|\mu'\| = \|\mu\|\) .
\item
  For every measure \(\mu\) on a locally compact space \(X\) and every homeomorphism \(\sigma\) of \(X\) onto itself, let \(\mu_{\sigma}\) be the measure \(f \mapsto \mu(f \circ \sigma)\) .
\end{enumerate}

\begin{enumerate}
\def\labelenumi{\alph{enumi})}
\item
  Show that \(|\mu_{\sigma}| = |\mu|_{\sigma}\) .
\item
  Equip \(\mathcal{M}(\mathrm{X};\mathbf{C})\) with the vague topology and denote by A the set of continuous endomorphisms of the locally convex space \(\mathcal{M}(\mathrm{X};\mathbf{C})\) ; equip A with the topology of uniform convergence in the vaguely bounded subsets of \(\mathcal{M}(\mathrm{X};\mathbf{C})\) . Let \(\Gamma\) be the group of homeomorphisms of X onto itself, equipped with the topology \(\mathcal{T}_{\beta}\) defined in GT, X, §3, No.~5. For every \(\sigma \in \Gamma\) , let \(A_{\sigma}\) be the mapping \(\mu \mapsto \mu_{\sigma}\) , which belongs to \(\mathcal{A}\) ; show that the mapping \(\sigma \mapsto A_{\sigma}\) of \(\Gamma\) into \(\mathcal{A}\) is continuous.
\end{enumerate}

\begin{enumerate}
\def\labelenumi{\arabic{enumi})}
\setcounter{enumi}{17}
\tightlist
\item
  Let \(X\) be a compact space, \(\mu\) a measure on \(X\) such that \(\mu(1) = \| \mu \|\) . Show that \(\mu\) is a positive measure. (If \(\mu_1 = \mathcal{R}(\mu)\) , \(\mu_2 = \mathcal{I}(\mu)\) , note first that \(\mu_1(1) = \| \mu_1 \|\) and show that \(\mu_1\) is a positive measure by making use of Ch. II, §2, No.~2, Prop. 4; then prove that \(\mu_2 = 0\) by considering \(\mu(1 + if)\) for a suitable function \(f \in \mathcal{C}(X; \mathbf{R})\) .)
\end{enumerate}

\setcounter{exercisegroup}{2}
\edef\CurrentExerciseGroup{\arabic{exercisegroup}}
